
\section*{Part IV. Applications of atomic physics \hfill 15 / 15}
\subsection*{1. Atomic clocks measure a stable transition frequency}
A clock locks a local oscillator to a narrow atomic transition.
The modern SI second is defined by the unperturbed ground-state
hyperfine transition frequency of \(^{133}\mathrm{Cs}\):
\[
\boxed{\nu_{\rm Cs}=9\,192\,631\,770\ \mathrm{Hz}.}
\]
Cs has configuration \([\mathrm{Xe}]6s\), electronic
\(J=1/2\), nuclear \(I=7/2\), hence \(F=3,4\).
The clock transition is
\[
\boxed{\ket{F=3,m_F=0}\ \longleftrightarrow\
\ket{F=4,m_F=0},}
\]
not \(F=3\to F=3\). The \(m_F=0\) choice strongly reduces
first-order Zeeman sensitivity near zero field.
Optical clocks exploit much higher carrier frequencies,
e.g. strontium clock transitions around \(429\) THz,
and narrow, often nominally forbidden, transitions.
Trap/lattice configurations help reduce motion and
frequency shifts.

\subsection*{2. Linewidth, interrogation time and stability}
The quality factor is \(Q=\nu_0/\Delta\nu\).
Long-lived excited states have narrow natural linewidths;
long coherent interrogation time \(T\) can produce spectral
features with width of order \(1/T\).
For \(N\) independent atoms, projection noise typically
scales as \(1/\sqrt N\) and repeated measurements improve
statistics as the square root of averaging time under
appropriate noise models. Real clock accuracy requires
controlling black-body shifts, Zeeman and Stark shifts,
collisions, gravitational redshift and laser noise.
Clock comparisons can test relativity and variation of
fundamental constants.

\subsection*{3. Find the right formula quickly}
\begin{center}\small
\begin{tabularx}{\linewidth}{@{}lXr@{}}\toprule
Question type & Starting formula or test & Page\\\midrule
Photon counts & \(p_n,\ \bar n,\ (\Delta n)^2,\ g^{(2)}\) & 3--4\\
Coherent/squeezed light & \(\hat a\ket\alpha=\alpha\ket\alpha,\ (\Delta x)^2=e^{-2r}\) & 2,4\\
Beam splitter / HOM & transform \(\hat a^\dagger,\hat b^\dagger\) & 3\\
Photonic entanglement & Bell basis; partial trace / entropy & 5\\
Teleport / QKD & measurement branches; basis sifting & 6\\
Optical loss & joint unitary then \(\operatorname{Tr}_E\) & 7\\
Photon subtraction & project detector mode, then normalise & 8\\
Cavity coupling & JC excitation block; dressed eigenvalues & 8\\
Interferometric sensing & \(F_Q=4\operatorname{Var}(G)\) & 9\\
Radial / finite nucleus & \(u=rR,\ \Delta E=\langle\Delta V\rangle\) & 10\\
Term symbols and \(3j\) & triangle rule, microstate count, W--E & 11\\
Hyperfine / Zeeman / Stark & \(E_F,\ g_Fm_FB,\ -\alpha E^2/2\) & 12\\
E1 / M1 / E2 and lifetimes & rank + parity + \(A_{if}\), then \(\tau\) & 13\\
Atomic PV / EDM & weak-induced mixing / \(-\mathbf d\cdot\mathbf E\) & 14\\
Clocks & \(F,m_F\) choice; \(\nu_0/\Delta\nu\) & 15\\\bottomrule
\end{tabularx}
\end{center}
\begin{methodbox}
\textbf{Universal exam workflow:} recognise the model and basis;
write assumptions/conventions; derive or quote the governing operator;
calculate with explicit normalisation; check dimensions,
limiting cases, parity/selection rules and degeneracy;
interpret the observable physically.
\end{methodbox}
\src{Week 13 pp.~1--35; Week 9 on hyperfine shifts; Week 10 on linewidths and forbidden lines.}
